% Academic CV template for Guillaume Huber
% Compile with pdfLaTeX on Overleaf.
% Standalone version: no external .cls file required.

\documentclass[11pt,letterpaper]{article}

\usepackage[margin=0.85in]{geometry}
\usepackage[T1]{fontenc}
\usepackage[utf8]{inputenc}
\usepackage{lmodern}
\usepackage{tabularx}
\usepackage{enumitem}
\usepackage{titlesec}
\usepackage{xcolor}
\usepackage[hidelinks]{hyperref}

\setlength{\parindent}{0pt}
\setlength{\parskip}{3pt}
\setlist[itemize]{leftmargin=*, topsep=2pt, itemsep=2pt, parsep=0pt}
\renewcommand{\arraystretch}{1.12}

\definecolor{cvblue}{RGB}{30,70,120}

\titleformat{\section}
  {\large\bfseries\color{cvblue}}
  {}{0em}{}[\titlerule]
\titlespacing*{\section}{0pt}{10pt}{5pt}

\newcommand{\cvheading}[4]{%
  \textbf{#1} \hfill #2\\
  \textit{#3} \hfill \textit{#4}\par
}

\newcommand{\cvsubheading}[2]{%
  \textbf{#1} \hfill #2\par
}

\newcommand{\publication}[4]{%
  \item #1. ``#2.'' \textit{#3}. #4.
}

\begin{document}

% =========================
% Header
% =========================

{\LARGE \textbf{Guillaume Huber}}\\[-1pt]
Graduate Student in Astronomy, Institute for Astronomy, University of Hawai`i at M\=anoa\\
Hilo, HI, USA \; | \; \href{mailto:ghuber@hawaii.edu}{ghuber@hawaii.edu} \; | \; +1 (631) 708--1965\\
% Optional: add website, ADS profile, ORCID, GitHub, LinkedIn
% \href{https://your-website.com}{your-website.com} \; | \; \href{https://orcid.org/XXXX-XXXX-XXXX-XXXX}{ORCID: XXXX-XXXX-XXXX-XXXX}








% =========================
% Research profile
% =========================

\section{Research Interests}
Astronomical instrumentation; UV-visible-NIR detector development; photon-counting and ultra-low-background detector systems; high-contrast imaging for exoplanets; optical and polarimetric instrumentation; precursor science and instrumentation for the Habitable Worlds Observatory.








% =========================
% Education
% =========================

\section{Education}

\cvheading{University of Hawai`i at M\=anoa, Institute for Astronomy}{2022--present}{Ph.D. student in Astronomy, M.Sc. in Astronomy}{Honolulu / Hilo, HI, USA}
\begin{itemize}
  \item Thesis topic: near-infrared arrays and target selection for the Habitable Worlds Observatory.
  \item Advisor: Prof. Michael Bottom
  \item Research areas: infrared detector characterization, astronomical instrumentation, polarimetry, high-contrast imaging, and target selection for future space missions.
\end{itemize}

\cvheading{Institut d'Optique Graduate School / Universit\'e Paris-Saclay}{2019--2022}{M.Eng. in Applied Physics, Optics and Photonics}{Palaiseau, France}
\begin{itemize}
  \item Coursework in optical instrumentation, Fourier optics, polarimetry, signal and image processing, quantum optics, radiometry, detection systems, lasers, optical fibers, and applied photonics.
  \item GPA: 4.36/4.5; graduated top 10\%.
\end{itemize}

\cvheading{Panthéon-Sorbonne University}{2021--2022}{First-year B.A. coursework in Philosophy}{Paris, France}
\begin{itemize}
  \item Coursework in general philosophy, philosophy of science, history of philosophy, and practical philosophy.
\end{itemize}

\cvheading{Classes Pr\'eparatoires aux Grandes \'Ecoles}{2017--2019}{B.Sc.-level intensive track in Mathematics and Physics}{Sophia Antipolis, Nantes, Rennes, France}
\begin{itemize}
  \item Highly selective preparatory program with intensive coursework in mathematics, physics, engineering, computer science, chemistry, and philosophy.
\end{itemize}








% =========================
% Research appointments
% =========================

\section{Research Experience}

\cvheading{Detector Characterization for Low-Signal Astronomy}{Aug. 2022--present}{Institute for Astronomy, University of Hawai`i - UC Berkeley}{Advisor: Prof. Michael Bottom}
\begin{itemize}
  \item Characterize large-format linear-mode avalanche photodiode (LmAPD) arrays for ultra-low-signal near-infrared astronomy and future astronomical observatories, including the Habitable Worlds Observatory.
  \item Develop laboratory experiments and analysis pipelines for detector characterization (gain, read noise, dark current, glow, persistence, and single-photon event detection).
  \item Analyze laboratory and telescope datasets; build simulations and statistical tools for detector performance assessment.
  \item Collaborate with academic partners (RIT, MIT, UC Berkeley), NASA, and industry partners (Leonardo UK, Hawaii Aerospace); presented related work at SPIE Astronomical Telescopes + Instrumentation, astronomy and HWO-related conferences, and detector workshops.
\end{itemize}

\cvheading{Target Selection for the Habitable Worlds Observatory}{Oct. 2024--present}{Institute for Astronomy, University of Hawai`i}{Advisors: Prof. Michael Bottom}
\begin{itemize}
  \item Search for binary companions in potential HWO target systems using Robo-AO-2 observations from the UH 88-inch telescope on Mauna Kea.
  \item Contribute to the instrument commissioning and its infrared data-processing pipeline
  \item Contribute to the HWO target-selection working group.
\end{itemize}

\cvheading{Dual-Wavelength Spectropolarimeter for Solar Astronomy}{Summer 2021; Aug. 2022--2024}{Institute for Astronomy, University of Hawai`i}{Advisor: Prof. Haosheng Lin}
\begin{itemize}
  \item Designed and demonstrated a single-channel dual-wavelength polarimetric concept using optical modeling, Mueller-matrix analysis, and laboratory validation.
  \item Built Python simulations and experimental performance tests for the polarimeter proof of concept.
  \item First-year project awarded the IfA Achievement Scholarship.
\end{itemize}

\cvheading{Starspots on Main-Sequence Superflaring Stars}{2023-2024}{Institute for Astronomy, University of Hawai`i}{Advisor: Prof. Haosheng Lin}
\begin{itemize}
  \item Measure stellar spot temperatures and magnetic-field diagnostics using TESS photometry and high-resolution infrared spectroscopy.
  \item Awarded observing time as Principal Investigator on NASA IRTF.
\end{itemize}

\cvheading{Optical Engineering for Greenhouse-Gas Earth Observation}{Mar.--June 2022}{SRON Netherlands Institute for Space Research}{Advisor: Dr. Ryan Cooney}
\begin{itemize}
  \item Contributed to an ESA-funded project for characterization of near-infrared cameras for the TANGO mission.
  \item Installed and evaluated a cleanroom sensor-testing setup involving cryogenic cooling and ultra-low-pressure conditions.
  \item Modeled thermal background and optical performance for infrared detector testing.
\end{itemize}

\cvheading{Hypertelescope Optical Design and Simulation}{July--Nov. 2021}{Hypertelescope LISE}{Advisors: Prof. Antoine Labeyrie; Prof. Thierry L\'epine}
\begin{itemize}
  \item Studied coma-aberration correction in a multi-pupil hypertelescope optical design.
  \item Developed Fourier-optics simulations in Python to model hypertelescope image formation.
\end{itemize}

% Optional: keep only if applying to instrumentation-heavy academic or engineering roles.
\cvheading{Junior Engineer Consulting}{2021--2022}{Valeo; Somfy}{France}
\begin{itemize}
  \item Characterized polymer automotive lenses, including focal length, surface geometry, thickness, and aberration measurements.
  \item Performed 3D photometric characterization and environmental performance tests of an embedded infrared detector system for home automation.
\end{itemize}








% =========================
% Publications
% =========================

\section{Publications}
% Replace placeholder author lists with full citation format when available.
% Recommended format: authors, year, title, journal/proceedings, volume, pages, DOI/arXiv/ADS link.

\subsection*{Published / Accepted}
\begin{enumerate}[leftmargin=*, label={[P\arabic*]}]
  \publication
{Huber, G. et al}
{Glow reduction of ultra-low noise LmAPDs: towards photon-counting infrared arrays}
{SPIE Journal of Astronomical Telescopes, Instruments, and Systems}
{2024. \href{https://ui.adsabs.harvard.edu/abs/2024SPIE13103E..20H/abstract}{ADS}}
\publication{Huber, G. et al}{Recent characterization results of LmAPDs: towards photon-counting infrared arrays for HWO}{HWO 2025 Conference Proceedings}{2026. \href{https://ui.adsabs.harvard.edu/abs/2026ASPC..543..345H/abstract}{ADS}}
  \publication{Huber, G., Lin, H}{Simultaneous multi-wavelength polarimetry without spectral splitting: proof of concept of a single-channel dual-wavelength polarimeter}{SPIE Journal of Astronomical Telescopes, Instruments, and Systems}{2024.  \href{https://ui.adsabs.harvard.edu/abs/2024SPIE13100E..3DH/abstract}{ADS}}
  \publication{Claveau, C.-A. et al}{Progress towards a megapixel linear-mode avalanche photodiode array for ultra-low-background shortwave infrared astronomy}{SPIE Journal of Astronomical Telescopes, Instruments, and Systems}{2024. \href{https://ui.adsabs.harvard.edu/abs/2024arXiv241109185C/abstract}{ADS}}
  \publication{Baranec, C. et al}{Robo-AO-2: entering the era of automated science operations, hybrid wavefront sensing, and adaptive secondary integration}{SPIE Journal of Astronomical Telescopes, Instruments, and Systems}{2026.  \href{https://ui.adsabs.harvard.edu/abs/2026arXiv260702826B/abstract}{ADS}}
  \publication{Hartman Z. et al}{Paving the Road to the Habitable Worlds Observatory with High-resolution Imaging. I. New and Archival Speckle Observations of Potential HWO Target Stars}{The Astronomical Journal}{2026. \href{https://ui.adsabs.harvard.edu/abs/2026AJ....171..174H/abstract}{ADS}}
  \publication{Marcussen M. et al}{The $\mu$ Herculis system solved after nearly three centuries}{Astronomy and Astrophysics}{2026. \href{https://doi.org/10.1051/0004-6361/202658981}{DOI}}
  \publication{Hartman Z. et al}{Here Be Binaries: Examining Our Current Understanding of the Multiplicity Status of Nearby Potential HWO Targets}{HWO 2025 Conference Proceedings}{Submitted, 2025}
\end{enumerate}

\subsection*{Submitted / In Preparation}
\begin{enumerate}[leftmargin=*, label={[S\arabic*]}]
  \publication{Huber, G. et al}{Near-noiseless single-photon detection in the infrared with a megapixel semiconductor array for low-background astronomy}{Nature Astronomy - Publications of the Astronomical Society of the Pacific}{In preparation}
  \publication{Huber, G. et al}{Persistence in Linear-mode Avalanche Photodiodes}{SPIE Journal of Astronomical Telescopes, Instruments, and Systems}{In preparation}
  \publication{Huber, G., Lin, H}{Proof of concept of a novel spectropolarimeter}{Publications of the Astronomical Society of the Pacific}{In preparation}
\end{enumerate}






% =========================
% Presentations
% =========================

\section{Conference Presentations and Workshops}
% Split into Invited Talks, Contributed Talks, and Posters if the list becomes long.
% Add exact titles, dates, locations, and whether each was a talk or poster.

\begin{itemize}
  \item SPIE Astronomical Telescopes + Instrumentation, 2024, Yokohama, Japan
  \item Caltech Sagan Workshop, 2024, Pasadena, CA, USA
  \item NASA Habitable Worlds Observatory Face-to-Face Meeting, 2024, Rochester, NY, USA
  \item Habitable Worlds Observatory Conference, 2025, Washington, DC, USA
  \item Scientific Detector Workshop, 2025, Canberra, Australia
  \item Spirit of Lyot Conference, 2026, Pasadena, CA, USA
\end{itemize}






% =========================
% Observing / proposals
% =========================

\section{Observing and Proposal Experience}
\begin{itemize}
    \item Principal Investigator, Robo-AO-2 observations of candidate HWO targets with the UH 88-inch telescope, Mauna Kea. 12 observing nights, survey of 508 targets.
  \item Principal Investigator, NASA IRTF/iSHELL observing program on spectroscopy of stellar spots on Main Sequence stars. 5 observing half-nights, 8 targets.
\end{itemize}






% =========================
% Teaching / mentoring
% =========================

\section{Teaching and Mentoring}
\begin{itemize}
  \item Private tutor in mathematics, physics, and English, 2019--2020.
  \item Elementary school lecture on galactic and extra-galactic astronomy for the \textit{Journey Through the Universe} outreach program in Hilo
\end{itemize}

% =========================
% Service / outreach
% =========================

\section{Service and Outreach}
\begin{itemize}
  \item Volunteering at IfA Open House, Astroday Hilo and Kona, and other outreach events
  \item Help setting up a 50cm Planewave telescope at Waipahu high-school
  \item Co-organizing \textit{Astronomy on Tap - Hawai`i}.
  \item Teaching at \textit{Journey through the Universe} in Hilo.
\end{itemize}

% =========================
% Awards
% =========================

\section{Awards and Scholarships}
\begin{itemize}
  \item Achievement Scholarship, Institute for Astronomy, University of Hawai`i, 2023. Merit award for the best first-year research project.
  \item Paris-Saclay University Travel Scholarship, 2021. Merit scholarship supporting travel and research stay in Hawai`i.
\end{itemize}

% =========================
% Technical skills
% =========================

\section{Technical Skills}
\begin{tabularx}{\textwidth}{@{}l X@{}}
Programming and data analysis & Python, MATLAB; basic IDL and C; scientific computing, statistical analysis, image processing, detector-data reduction.\\
Optics and instrumentation & Optical alignment, cryogenic detector testing, polarimetry, Fourier optics, radiometry, laboratory automation, cleanroom sensor testing.\\
Software & Zemax OpticStudio, basic SolidWorks, LaTeX, DS9, astropy, photutils, emcee\\
Languages & French native; English fluent; German intermediate; Japanese beginner.\\
\end{tabularx}

% =========================
% Optional professional experience
% =========================

\section{Selected Professional Experience}
% Keep this section short in an academic CV unless applying to instrumentation or engineering-heavy roles.

\cvsubheading{Co-founder and CTO, LensCare}{2019--2020}
\begin{itemize}
  \item Led technical research, prototyping, and early production for a student startup developing ecological contact-lens cases using UV-visible diode illumination.
\end{itemize}

% =========================
% References
% =========================

\section{References}
Available upon request.

% Alternative for academic applications:
% \begin{tabularx}{\textwidth}{@{}X X@{}}
% Prof. Michael Bottom & Prof. Haosheng Lin\\
% Institute for Astronomy, University of Hawai`i & Institute for Astronomy, University of Hawai`i\\
% \href{mailto:...}{...} & \href{mailto:...}{...}\\
% \end{tabularx}

\end{document}
